% Options for packages loaded elsewhere
\PassOptionsToPackage{unicode}{hyperref}
\PassOptionsToPackage{hyphens}{url}
\PassOptionsToPackage{dvipsnames,svgnames,x11names}{xcolor}
\documentclass[
  spanish,
  10pt,
  a4paper,
]{article}
\usepackage{xcolor}
\usepackage[top=2.2cm,bottom=2.2cm,left=2.1cm,right=2.1cm]{geometry}
\usepackage{amsmath,amssymb}
\setcounter{secnumdepth}{5}
\usepackage{iftex}
\ifPDFTeX
  \usepackage[T1]{fontenc}
  \usepackage[utf8]{inputenc}
  \usepackage{textcomp} % provide euro and other symbols
\else % if luatex or xetex
  \usepackage{unicode-math} % this also loads fontspec
  \defaultfontfeatures{Scale=MatchLowercase}
  \defaultfontfeatures[\rmfamily]{Ligatures=TeX,Scale=1}
\fi
\usepackage{lmodern}
\ifPDFTeX\else
  % xetex/luatex font selection
\fi
% Use upquote if available, for straight quotes in verbatim environments
\IfFileExists{upquote.sty}{\usepackage{upquote}}{}
\IfFileExists{microtype.sty}{% use microtype if available
  \usepackage[]{microtype}
  \UseMicrotypeSet[protrusion]{basicmath} % disable protrusion for tt fonts
}{}
\makeatletter
\@ifundefined{KOMAClassName}{% if non-KOMA class
  \IfFileExists{parskip.sty}{%
    \usepackage{parskip}
  }{% else
    \setlength{\parindent}{0pt}
    \setlength{\parskip}{6pt plus 2pt minus 1pt}}
}{% if KOMA class
  \KOMAoptions{parskip=half}}
\makeatother
\usepackage{color}
\usepackage{fancyvrb}
\newcommand{\VerbBar}{|}
\newcommand{\VERB}{\Verb[commandchars=\\\{\}]}
\DefineVerbatimEnvironment{Highlighting}{Verbatim}{commandchars=\\\{\}}
% Add ',fontsize=\small' for more characters per line
\newenvironment{Shaded}{}{}
\newcommand{\AlertTok}[1]{\textcolor[rgb]{1.00,0.00,0.00}{\textbf{#1}}}
\newcommand{\AnnotationTok}[1]{\textcolor[rgb]{0.38,0.63,0.69}{\textbf{\textit{#1}}}}
\newcommand{\AttributeTok}[1]{\textcolor[rgb]{0.49,0.56,0.16}{#1}}
\newcommand{\BaseNTok}[1]{\textcolor[rgb]{0.25,0.63,0.44}{#1}}
\newcommand{\BuiltInTok}[1]{\textcolor[rgb]{0.00,0.50,0.00}{#1}}
\newcommand{\CharTok}[1]{\textcolor[rgb]{0.25,0.44,0.63}{#1}}
\newcommand{\CommentTok}[1]{\textcolor[rgb]{0.38,0.63,0.69}{\textit{#1}}}
\newcommand{\CommentVarTok}[1]{\textcolor[rgb]{0.38,0.63,0.69}{\textbf{\textit{#1}}}}
\newcommand{\ConstantTok}[1]{\textcolor[rgb]{0.53,0.00,0.00}{#1}}
\newcommand{\ControlFlowTok}[1]{\textcolor[rgb]{0.00,0.44,0.13}{\textbf{#1}}}
\newcommand{\DataTypeTok}[1]{\textcolor[rgb]{0.56,0.13,0.00}{#1}}
\newcommand{\DecValTok}[1]{\textcolor[rgb]{0.25,0.63,0.44}{#1}}
\newcommand{\DocumentationTok}[1]{\textcolor[rgb]{0.73,0.13,0.13}{\textit{#1}}}
\newcommand{\ErrorTok}[1]{\textcolor[rgb]{1.00,0.00,0.00}{\textbf{#1}}}
\newcommand{\ExtensionTok}[1]{#1}
\newcommand{\FloatTok}[1]{\textcolor[rgb]{0.25,0.63,0.44}{#1}}
\newcommand{\FunctionTok}[1]{\textcolor[rgb]{0.02,0.16,0.49}{#1}}
\newcommand{\ImportTok}[1]{\textcolor[rgb]{0.00,0.50,0.00}{\textbf{#1}}}
\newcommand{\InformationTok}[1]{\textcolor[rgb]{0.38,0.63,0.69}{\textbf{\textit{#1}}}}
\newcommand{\KeywordTok}[1]{\textcolor[rgb]{0.00,0.44,0.13}{\textbf{#1}}}
\newcommand{\NormalTok}[1]{#1}
\newcommand{\OperatorTok}[1]{\textcolor[rgb]{0.40,0.40,0.40}{#1}}
\newcommand{\OtherTok}[1]{\textcolor[rgb]{0.00,0.44,0.13}{#1}}
\newcommand{\PreprocessorTok}[1]{\textcolor[rgb]{0.74,0.48,0.00}{#1}}
\newcommand{\RegionMarkerTok}[1]{#1}
\newcommand{\SpecialCharTok}[1]{\textcolor[rgb]{0.25,0.44,0.63}{#1}}
\newcommand{\SpecialStringTok}[1]{\textcolor[rgb]{0.73,0.40,0.53}{#1}}
\newcommand{\StringTok}[1]{\textcolor[rgb]{0.25,0.44,0.63}{#1}}
\newcommand{\VariableTok}[1]{\textcolor[rgb]{0.10,0.09,0.49}{#1}}
\newcommand{\VerbatimStringTok}[1]{\textcolor[rgb]{0.25,0.44,0.63}{#1}}
\newcommand{\WarningTok}[1]{\textcolor[rgb]{0.38,0.63,0.69}{\textbf{\textit{#1}}}}
\usepackage{longtable,booktabs,array}
\newcounter{none} % for unnumbered tables
\usepackage{calc} % for calculating minipage widths
% Correct order of tables after \paragraph or \subparagraph
\usepackage{etoolbox}
\makeatletter
\patchcmd\longtable{\par}{\if@noskipsec\mbox{}\fi\par}{}{}
\makeatother
% Allow footnotes in longtable head/foot
\IfFileExists{footnotehyper.sty}{\usepackage{footnotehyper}}{\usepackage{footnote}}
\makesavenoteenv{longtable}
\ifLuaTeX
\usepackage[bidi=basic,shorthands=off]{babel}
\else
\usepackage[bidi=default,shorthands=off]{babel}
\fi
\ifLuaTeX
  \usepackage{selnolig} % disable illegal ligatures
\fi
\setlength{\emergencystretch}{3em} % prevent overfull lines
\providecommand{\tightlist}{%
  \setlength{\itemsep}{0pt}\setlength{\parskip}{0pt}}
% Shared visual layer for the generated Kaleido FlowTwin reports.
\usepackage{helvet}
\renewcommand{\familydefault}{\sfdefault}
\usepackage{microtype}
\usepackage{xcolor}
\usepackage{titlesec}
\usepackage{fancyhdr}
\usepackage{enumitem}
\usepackage{booktabs}
\usepackage{xurl}
\usepackage{lastpage}
\usepackage{etoolbox}

\definecolor{KaleidoNavy}{HTML}{0A2342}
\definecolor{KaleidoBlue}{HTML}{176B87}
\definecolor{KaleidoCyan}{HTML}{3BB7C4}
\definecolor{KaleidoMint}{HTML}{DDF5F1}
\definecolor{KaleidoFog}{HTML}{F2F5F7}
\definecolor{KaleidoSlate}{HTML}{5D6A7A}
\definecolor{KaleidoCoral}{HTML}{F26B5B}

\titleformat{\section}
  {\Large\bfseries\color{KaleidoNavy}}
  {\thesection}{0.75em}{}[\vspace{0.2em}\color{KaleidoCyan}\titlerule]
\titleformat{\subsection}
  {\large\bfseries\color{KaleidoBlue}}
  {\thesubsection}{0.65em}{}
\titleformat{\subsubsection}
  {\normalsize\bfseries\color{KaleidoNavy}}
  {\thesubsubsection}{0.55em}{}
\titlespacing*{\section}{0pt}{2.2em}{1em}
\titlespacing*{\subsection}{0pt}{1.6em}{0.65em}

\pagestyle{fancy}
\fancyhf{}
\fancyhead[L]{\small\bfseries\color{KaleidoNavy}KALEIDO FLOWTWIN}
\fancyhead[R]{\small\color{KaleidoSlate}DOSSIER TECNICO}
\fancyfoot[L]{\scriptsize\color{KaleidoSlate}EVOCON Solutions GmbH}
\fancyfoot[R]{\scriptsize\color{KaleidoSlate}\thepage/\pageref*{LastPage}}
\renewcommand{\headrulewidth}{0.6pt}
\renewcommand{\headrule}{\hbox to\headwidth{\color{KaleidoCyan}\leaders\hrule height \headrulewidth\hfill}}
\renewcommand{\footrulewidth}{0pt}
\setlength{\headheight}{20pt}

\setlist[itemize]{leftmargin=1.4em,itemsep=0.25em,topsep=0.35em}
\setlist[enumerate]{leftmargin=1.7em,itemsep=0.25em,topsep=0.35em}
\setlength{\parindent}{0pt}
\setlength{\parskip}{0.55em}
\setlength{\emergencystretch}{3em}
\renewcommand{\arraystretch}{1.18}

\AtBeginEnvironment{quote}{%
  \color{KaleidoNavy}\itshape
  \setlength{\leftskip}{1em}
  \setlength{\rightskip}{1em}
}

\makeatletter
\renewcommand{\maketitle}{%
  \hypersetup{pageanchor=false}
  \begin{titlepage}
    \pagecolor{KaleidoNavy}
    \color{white}
    \vspace*{1.4cm}
    {\color{KaleidoCyan}\rule{2.4cm}{2.5pt}\par}
    \vspace{1.1cm}
    {\Huge\bfseries\@title\par}
    \vspace{0.65cm}
    {\Large\color{KaleidoMint}Kaleido FlowTwin\par}
    \vfill
    {\large Álvaro Schwiedop Souto\par}
    \vspace{0.15cm}
    {\normalsize\color{KaleidoMint}Industrial AI \& Predictive Quality Lead\par}
    \vspace{0.15cm}
    {\normalsize\color{KaleidoMint}EVOCON Solutions GmbH\par}
    \vspace{0.25cm}
    {\normalsize\color{KaleidoMint}\@date\par}
  \end{titlepage}
  \hypersetup{pageanchor=true}
  \nopagecolor
  \color{black}
}
\makeatother
\usepackage{bookmark}
\IfFileExists{xurl.sty}{\usepackage{xurl}}{} % add URL line breaks if available
\urlstyle{same}
% fallback for those not using the hyperref driver hyperxmp:
\makeatletter
\@ifundefined{xmpquote}{\newcommand{\xmpquote}[1]{#1}}{}
\makeatother
\hypersetup{
  pdftitle={Arquitectura de Kaleido FlowTwin},
  pdflang={es},
  colorlinks=true,
  linkcolor={KaleidoBlue},
  filecolor={Maroon},
  citecolor={Blue},
  urlcolor={KaleidoBlue},
  pdfcreator={LaTeX via pandoc}}

\title{Arquitectura de Kaleido FlowTwin}
\author{}
\date{15 julio 2026}

\begin{document}
\maketitle

{
\setcounter{tocdepth}{3}
\tableofcontents
}
\section{Arquitectura de Kaleido
FlowTwin}\label{arquitectura-de-kaleido-flowtwin}

\subsection{Principios}\label{principios}

\begin{enumerate}
\def\labelenumi{\arabic{enumi}.}
\tightlist
\item
  \textbf{Data-poor by design.} El sistema entrega auditoria y process
  mining antes de necesitar un modelo neuronal.
\item
  \textbf{Multiobjeto.} Proyecto, operacion, turno, unidad de carga,
  recurso, buque, documento e incidencia coexisten; no se fuerza todo a
  un unico case ID.
\item
  \textbf{Plan versionado.} Plan inicial y replans son eventos, no
  columnas sobreescritas.
\item
  \textbf{Accion no es contexto.} Solo una decision timestamped e
  intervenible es accion.
\item
  \textbf{Read-only.} El MVP asesora; no controla equipos ni sistemas.
\item
  \textbf{Baselines first.} Un modelo complejo se conserva solo si anade
  valor.
\item
  \textbf{Uncertainty first.} Toda prediccion incluye intervalo,
  confianza o abstencion.
\end{enumerate}

\subsection{Vista logica}\label{vista-logica}

\begin{Shaded}
\begin{Highlighting}[]
\NormalTok{Trace Port / Shipping Board / Freight Intelligence / CSV}
\NormalTok{                          |}
\NormalTok{                  read{-}only extract}
\NormalTok{                          v}
\NormalTok{          Canonical Operational Event Contract}
\NormalTok{         ids {-} event time {-} plan time {-} roles {-} units}
\NormalTok{                          |}
\NormalTok{              validation + lineage + quality}
\NormalTok{                          |}
\NormalTok{              object{-}centric event graph}
\NormalTok{       project{-}operation{-}shift{-}cargo{-}resource{-}vessel}
\NormalTok{                          |}
\NormalTok{        +{-}{-}{-}{-}{-}{-}{-}{-}{-}{-}{-}{-}{-}{-}{-}{-}{-}+{-}{-}{-}{-}{-}{-}{-}{-}{-}{-}{-}{-}{-}{-}{-}{-}{-}{-}+}
\NormalTok{        |                 |                  |}
\NormalTok{ process mining     predictive baselines   simulation}
\NormalTok{ conformance        boosting/survival      discrete{-}event}
\NormalTok{        |                 |                  |}
\NormalTok{        +{-}{-}{-}{-}{-}{-}{-}{-}{-}{-}{-}{-}{-}{-}{-}{-}{-}+{-}{-}{-}{-}{-}{-}{-}{-}{-}{-}{-}{-}{-}{-}{-}{-}{-}{-}+}
\NormalTok{                          |}
\NormalTok{              optional Event{-}JEPA module}
\NormalTok{              future latent state by horizon}
\NormalTok{                          |}
\NormalTok{             calibration + value evaluation}
\NormalTok{                          |}
\NormalTok{       dashboard / API / CSV / PDF / audit trail}
\end{Highlighting}
\end{Shaded}

\subsection{Contrato canonico}\label{contrato-canonico}

\subsubsection{\texorpdfstring{\texttt{OperationEvent}}{OperationEvent}}\label{operationevent}

\begin{Shaded}
\begin{Highlighting}[]
\AttributeTok{@dataclass}\NormalTok{(frozen}\OperatorTok{=}\VariableTok{True}\NormalTok{)}
\KeywordTok{class}\NormalTok{ OperationEvent:}
\NormalTok{    event\_id: }\BuiltInTok{str}
\NormalTok{    source\_system: }\BuiltInTok{str}
\NormalTok{    source\_record\_id: }\BuiltInTok{str}
\NormalTok{    event\_type: }\BuiltInTok{str}
\NormalTok{    event\_time\_utc: datetime}
\NormalTok{    ingested\_at\_utc: datetime}
\NormalTok{    project\_id: }\BuiltInTok{str} \OperatorTok{|} \VariableTok{None}
\NormalTok{    operation\_id: }\BuiltInTok{str} \OperatorTok{|} \VariableTok{None}
\NormalTok{    shift\_id: }\BuiltInTok{str} \OperatorTok{|} \VariableTok{None}
\NormalTok{    cargo\_unit\_id: }\BuiltInTok{str} \OperatorTok{|} \VariableTok{None}
\NormalTok{    resource\_id: }\BuiltInTok{str} \OperatorTok{|} \VariableTok{None}
\NormalTok{    vessel\_id: }\BuiltInTok{str} \OperatorTok{|} \VariableTok{None}
\NormalTok{    location\_id: }\BuiltInTok{str} \OperatorTok{|} \VariableTok{None}
\NormalTok{    actor\_role: }\BuiltInTok{str} \OperatorTok{|} \VariableTok{None}
\NormalTok{    status\_from: }\BuiltInTok{str} \OperatorTok{|} \VariableTok{None}
\NormalTok{    status\_to: }\BuiltInTok{str} \OperatorTok{|} \VariableTok{None}
\NormalTok{    numeric\_value: }\BuiltInTok{float} \OperatorTok{|} \VariableTok{None}
\NormalTok{    unit: }\BuiltInTok{str} \OperatorTok{|} \VariableTok{None}
\NormalTok{    data\_quality\_flags: }\BuiltInTok{tuple}\NormalTok{[}\BuiltInTok{str}\NormalTok{, ...]}
\NormalTok{    payload\_ref: }\BuiltInTok{str} \OperatorTok{|} \VariableTok{None}
\end{Highlighting}
\end{Shaded}

\subsubsection{Plan y resultado}\label{plan-y-resultado}

\begin{Shaded}
\begin{Highlighting}[]
\AttributeTok{@dataclass}\NormalTok{(frozen}\OperatorTok{=}\VariableTok{True}\NormalTok{)}
\KeywordTok{class}\NormalTok{ PlanRevision:}
\NormalTok{    plan\_id: }\BuiltInTok{str}
\NormalTok{    revision: }\BuiltInTok{int}
\NormalTok{    valid\_from\_utc: datetime}
\NormalTok{    operation\_id: }\BuiltInTok{str}
\NormalTok{    milestone: }\BuiltInTok{str}
\NormalTok{    planned\_start\_utc: datetime }\OperatorTok{|} \VariableTok{None}
\NormalTok{    planned\_end\_utc: datetime }\OperatorTok{|} \VariableTok{None}
\NormalTok{    planned\_quantity: }\BuiltInTok{float} \OperatorTok{|} \VariableTok{None}
\NormalTok{    planned\_resources: }\BuiltInTok{tuple}\NormalTok{[}\BuiltInTok{str}\NormalTok{, ...]}
\NormalTok{    reason: }\BuiltInTok{str} \OperatorTok{|} \VariableTok{None}

\AttributeTok{@dataclass}\NormalTok{(frozen}\OperatorTok{=}\VariableTok{True}\NormalTok{)}
\KeywordTok{class}\NormalTok{ OperationOutcome:}
\NormalTok{    operation\_id: }\BuiltInTok{str}
\NormalTok{    completed\_at\_utc: datetime }\OperatorTok{|} \VariableTok{None}
\NormalTok{    completion\_status: }\BuiltInTok{str}
\NormalTok{    deviation\_minutes: }\BuiltInTok{float} \OperatorTok{|} \VariableTok{None}
\NormalTok{    incident\_types: }\BuiltInTok{tuple}\NormalTok{[}\BuiltInTok{str}\NormalTok{, ...]}
\NormalTok{    cost\_eur: }\BuiltInTok{float} \OperatorTok{|} \VariableTok{None}
\NormalTok{    censored: }\BuiltInTok{bool}
\end{Highlighting}
\end{Shaded}

Invariantes:

\begin{itemize}
\tightlist
\item
  timestamps originales no se reescriben;
\item
  toda correccion crea nueva revision;
\item
  unidades se conservan y normalizan con trazabilidad;
\item
  IDs se seudonimizan de forma estable;
\item
  outcome futuro nunca entra como feature;
\item
  una foto es referencia a objeto, no blob dentro del log;
\item
  missing es explicito, no cero;
\item
  eventos duplicados se marcan, no se borran silenciosamente.
\end{itemize}

\subsection{Roles de columnas}\label{roles-de-columnas}

\subsubsection{Acciones validas}\label{acciones-validas}

\begin{itemize}
\tightlist
\item
  asignar/cambiar equipo o recurso;
\item
  cambiar secuencia de carga/descarga;
\item
  aprobar pausa/reinicio;
\item
  cambiar ventana o turno;
\item
  reasignar posicion/area;
\item
  solicitar inspeccion adicional;
\item
  actualizar plan de forma timestamped.
\end{itemize}

\subsubsection{Contexto}\label{contexto}

\begin{itemize}
\tightlist
\item
  cliente, terminal, buque y muelle;
\item
  tipo, peso y dimensiones de carga;
\item
  turno/dia de semana;
\item
  meteo y estado de mar;
\item
  equipo disponible si no consta una decision;
\item
  procedimiento, puerto y pais;
\item
  complejidad de documentacion.
\end{itemize}

\subsubsection{Observaciones}\label{observaciones}

\begin{itemize}
\tightlist
\item
  milestones realizados;
\item
  cantidad movida acumulada;
\item
  tiempos entre eventos;
\item
  colas, esperas y pausas;
\item
  incidencias, notas y fotos;
\item
  estado de recursos;
\item
  posicion AIS o de unidad cuando exista.
\end{itemize}

\subsubsection{Outcomes}\label{outcomes}

\begin{itemize}
\tightlist
\item
  fin real;
\item
  desviacion respecto al plan congelado;
\item
  incidencia posterior;
\item
  dano/reclamacion;
\item
  coste y penalizacion;
\item
  SLA incumplido.
\end{itemize}

\subsection{Object-centric event
graph}\label{object-centric-event-graph}

Un evento puede relacionarse con varios objetos. Ejemplo:

\begin{Shaded}
\begin{Highlighting}[]
\NormalTok{event: CRANE\_LIFT\_COMPLETED}
\NormalTok{objects:}
\NormalTok{  operation=OP{-}42}
\NormalTok{  shift=SH{-}3}
\NormalTok{  cargo\_unit=CU{-}117}
\NormalTok{  resource=CRANE{-}2}
\NormalTok{  vessel=IMO{-}...}
\end{Highlighting}
\end{Shaded}

Relaciones principales:

\begin{Shaded}
\begin{Highlighting}[]
\NormalTok{project contains operation}
\NormalTok{operation scheduled\_in shift}
\NormalTok{operation handles cargo\_unit}
\NormalTok{resource executes event}
\NormalTok{vessel hosts operation}
\NormalTok{incident affects cargo\_unit/resource/operation}
\NormalTok{plan\_revision governs operation after valid\_from}
\end{Highlighting}
\end{Shaded}

El grafo evita duplicar eventos en varias trazas y permite explicar
propagacion.

\subsection{Escalera de modelos}\label{escalera-de-modelos}

\subsubsection{M0 - Descriptivo}\label{m0---descriptivo}

\begin{itemize}
\tightlist
\item
  calidad/cobertura del log;
\item
  tiempos por actividad;
\item
  variantes y directly-follows graph;
\item
  esperas y rework;
\item
  conformance respecto al flujo acordado.
\end{itemize}

\subsubsection{M1 - Baselines}\label{m1---baselines}

\begin{itemize}
\tightlist
\item
  mediana por clase de operacion;
\item
  Kaplan-Meier/Cox si hay censura;
\item
  quantile Gradient Boosting/CatBoost;
\item
  reglas de umbral y control charts;
\item
  regresion/logistica regularizada.
\end{itemize}

\subsubsection{M2 - Secuencial}\label{m2---secuencial}

\begin{itemize}
\tightlist
\item
  GRU/TCN;
\item
  ProcessTransformer;
\item
  temporal graph model object-centric;
\item
  PGTNet-style remaining-time baseline.
\end{itemize}

\subsubsection{M3 - Event-JEPA}\label{m3---event-jepa}

Estado:

\begin{Shaded}
\begin{Highlighting}[]
\NormalTok{z\_t = Encoder(prefix\_event\_graph\_t)}
\NormalTok{z\_hat\_t+h = Predictor(z\_t, action\_prefix, context, horizon)}
\end{Highlighting}
\end{Shaded}

Target:

\begin{Shaded}
\begin{Highlighting}[]
\NormalTok{z\_t+h = TargetEncoder(prefix\_event\_graph\_t+h)}
\end{Highlighting}
\end{Shaded}

Heads:

\begin{itemize}
\tightlist
\item
  remaining-time quantiles;
\item
  deviation risk;
\item
  next-event distribution;
\item
  incident risk;
\item
  bottleneck/object attribution.
\end{itemize}

\subsubsection{M4 - Escenarios}\label{m4---escenarios}

Dos rutas separadas:

\begin{enumerate}
\def\labelenumi{\arabic{enumi}.}
\tightlist
\item
  simulacion discreta calibrada para colas/recursos;
\item
  ranking con world model de acciones observadas.
\end{enumerate}

La primera funciona con conocimiento experto y poco dato. La segunda
exige acciones reales variadas, soporte y validacion
\texttt{correct\ vs\ shuffled}.

\subsection{Cold start}\label{cold-start}

{\def\LTcaptype{none} % do not increment counter
\begin{longtable}[]{@{}ll@{}}
\toprule\noalign{}
Datos disponibles & Entregable permitido \\
\midrule\noalign{}
\endhead
\bottomrule\noalign{}
\endlastfoot
0-5 operaciones & schema, instrumentacion, mapa de proceso, simulacion
experta \\
5-30 & distribuciones descriptivas, controles, retrieval de casos
similares \\
30-100 & boosting/quantiles con intervalos amplios y validacion grouped
CV \\
100-500 & modelos secuenciales pequenos, calibracion por grupo \\
\textgreater500 y diversidad suficiente & Event-JEPA/graph, ablations y
adaptacion \\
\end{longtable}
}

Los umbrales son orientativos; la diversidad de estados/eventos importa
mas que el numero bruto.

\subsection{Split y leakage}\label{split-y-leakage}

Prioridad:

\begin{enumerate}
\def\labelenumi{\arabic{enumi}.}
\tightlist
\item
  holdout cronologico futuro;
\item
  holdout por proyecto/cliente/puerto;
\item
  holdout por tipo de carga;
\item
  grouped cross-validation.
\end{enumerate}

Prohibido:

\begin{itemize}
\tightlist
\item
  dividir eventos de la misma operacion entre train/test;
\item
  usar la revision final del plan como si se conociera al inicio;
\item
  calcular features con eventos posteriores al prediction time;
\item
  seleccionar umbral con test;
\item
  usar fotos/informes generados despues del outcome;
\item
  convertir tiempo-a-fin en feature.
\end{itemize}

\subsection{Evaluacion}\label{evaluacion}

\subsubsection{Prediccion}\label{prediccion}

\begin{itemize}
\tightlist
\item
  MAE/MedAE de tiempo restante;
\item
  pinball loss P50/P90;
\item
  cobertura y anchura de intervalos;
\item
  AUPRC de desviacion/incidencia;
\item
  Brier/ECE;
\item
  next-event macro-F1/top-k.
\end{itemize}

\subsubsection{Operacion}\label{operacion}

\begin{itemize}
\tightlist
\item
  lead time antes de desviacion;
\item
  falsas alertas por turno/100 operaciones;
\item
  eventos no detectados;
\item
  precision top-10\% de operaciones;
\item
  utilidad de explicacion valorada por operador;
\item
  latencia y disponibilidad.
\end{itemize}

\subsubsection{Negocio}\label{negocio}

\begin{itemize}
\tightlist
\item
  horas de espera evitables identificadas;
\item
  overtime/demurrage potencial;
\item
  tiempo de reporting/coordinacion;
\item
  cumplimiento de SLA;
\item
  ahorro simulado bajo politica, separado de ahorro realizado.
\end{itemize}

\subsection{Serving}\label{serving}

Endpoints read-only candidatos:

\begin{Shaded}
\begin{Highlighting}[]
\NormalTok{POST /v1/events/batch}
\NormalTok{POST /v1/score/operation{-}prefix}
\NormalTok{GET  /v1/operations/\{id\}/risk}
\NormalTok{POST /v1/scenarios/rank}
\NormalTok{GET  /v1/models/\{version\}/card}
\NormalTok{GET  /health}
\end{Highlighting}
\end{Shaded}

Cada respuesta incluye:

\begin{Shaded}
\begin{Highlighting}[]
\NormalTok{model\_version}
\NormalTok{data\_cutoff}
\NormalTok{plan\_revision}
\NormalTok{prediction\_time}
\NormalTok{horizon}
\NormalTok{point\_estimate}
\NormalTok{interval}
\NormalTok{confidence}
\NormalTok{abstained}
\NormalTok{reason\_codes}
\NormalTok{source\_event\_ids}
\end{Highlighting}
\end{Shaded}

\subsection{Seguridad y privacidad}\label{seguridad-y-privacidad}

\begin{itemize}
\tightlist
\item
  export read-only y minimo privilegio;
\item
  seudonimizacion antes de salir del entorno acordado;
\item
  cifrado en transito/reposo;
\item
  sin credenciales en artefactos;
\item
  retencion y borrado acordados;
\item
  fotos y notas tratadas como potencial dato personal/confidencial;
\item
  auditoria por prediccion;
\item
  no usar datos de Kaleido para modelos de otros clientes sin acuerdo.
\end{itemize}

\subsection{Despliegue}\label{despliegue}

\begin{enumerate}
\def\labelenumi{\arabic{enumi}.}
\tightlist
\item
  offline replay;
\item
  shadow dashboard;
\item
  alertas no accionables visibles a equipo tecnico;
\item
  revision por operaciones;
\item
  alertas operativas con runbook;
\item
  escenarios asesorados;
\item
  integracion mas profunda solo tras aprobacion y seguridad.
\end{enumerate}

\end{document}
